% ============================================================
% Chapter 13 — Conclusion
% ============================================================
\chapter{Conclusion}
\label{ch:conclusion}

\roadmap{This chapter restates the problem, summarizes what was built and
learned, and states the project's contribution in a form suitable for the
examiner's final reading.}

This project set out to make the invisible visible: to render learners'
affective and cognitive states measurable during authentic online
learning, and to do so on the hardware every student already owns. The
motivation was the affect-dynamics account of learning
\citep{dmello2012dynamics, baker2010better}, in which the transitions
among engagement, confusion, frustration, and boredom shape outcomes yet
vanish from view the moment instruction moves online. The response was a
browser-based intelligent tutoring platform in which retrieval-grounded
LLM tutoring, four evidence-based strategy interventions, and a
multimodal sensing layer coexist, and in which every learning session
becomes a time-aligned, inspectable, re-codable multimodal record.

What was delivered exceeds the proposal in some dimensions and reframes it
in others, and the report has been candid about which is which
(\cref{ch:scope}). The authoring and tutoring stack met or exceeded its
promises; the sensing and analysis instrumentation---session replay, the
AOI/SEEV attention layer, executive-function text mining, retrospective
coding, and the exporters---was built largely beyond the original scope
and constitutes the project's principal contribution. Four proposal items
were reframed rather than delivered as written, and the report reports
them as partial or reframed without upgrade. The measurement methodology
was designed with equal candor: it validates one sensing channel against
another as a convergent, non-criterion standard
\citep{campbell1959}, anchors that comparison in periodic self-report, and
states plainly the two gaps---no stored phase marker and no stored gain
fields---that the statistics must work around.

The study's quantitative findings \todo{summarize in one or two sentences
once results are frozen: the observed detection agreement (RQ1) and the
observed learning gain with effect size (RQ2), each with its honest
qualifier}. Beyond those numbers, the enduring result is the instrument
itself. By keeping sensing decoupled from pedagogy, the platform provides
a clean observational base on which detection validity can be established
before any closed loop is attempted; by capturing dense, privacy-conscious,
time-aligned multimodal data with the tools to replay and code it, it
turns the measurement of learning into a tractable engineering problem.
The path from here---closing the affect-to-support loop on a validated
detector (\cref{ch:future})---is now a matter of building on a foundation
that this project has laid and, as far as the code allows one to say,
laid honestly.
